\section{Related work}

\paragraph{Interventions, steering, and interpretability benchmarks.}
Activation addition and representation engineering edit internal activations,
while feedback steering and local optimal-control methods place those edits in a
controller
\cite{turner2023activation,zou2023representation,nguyen2025feedback,skifstad2026local}.
Activation and attribution patching estimate component effects; AtP* provides a
local approximation; causal abstraction tests assigned causal roles; and pyvene
executes model interventions
\cite{syed2024attribution,kramar2024atp,geiger2024alignments,wu2024pyvene}.
Interactions between mediators can hide or inflate patching effects
\cite{vaidyanathan2026curse}. MIB, AxBench, and InterpBench instead benchmark
circuit localization, concept detection and steering, or recovery of known
structure \cite{mueller2025mib,wu2025axbench,gupta2024interpbench}. Mechanistic
tomography treats measurements as a way to recover a mechanism or effect map
\cite{erramilli2026tomography}; ObserverBench starts with such an estimate and
asks whether it supports a declared intervention or controller.

\paragraph{Decisions and AI control.}
Decision-focused learning evaluates a predictor through the downstream decision
rather than assuming that lower prediction error is enough
\cite{donti2017task,elmachtoub2022smart}. The same distinction motivates calls
to judge interpretability by the actions it enables \cite{orgad2026actionable}.
AI control asks whether a protocol can remain safe when an untrusted model may
try to evade it, and ControlArena provides reusable environments, policies,
monitors, and protocol-level evaluation
\cite{greenblatt2024aicontrol,inglis2025controlarena}. ObserverBench focuses on
the estimate inside the protocol: what it predicts, what it can read, which
actions it can affect, and the loss produced by those actions.

\paragraph{Borrowed testbeds.}
The IOI circuit supplies known structure rather than a new mechanism claim
\cite{wang2023ioi}. Conditional backup, Hydra-style self-repair, copy
suppression, and prompt-dependent self-repair are already documented
\cite{mcgrath2023hydra,mcdougall2023copy,rushing2024selfrepair}. Induction heads
support the repeated-key continuation used in the Qwen task
\cite{olsson2022induction}, and Qwen2.5-7B supplies a second model surface
\cite{qwen2024qwen25}. We use these systems to define interventions, then judge
observers on fresh held-out masks and prompts rather than claim a new IOI or
induction mechanism.
